% AI-written, AI cross-reviewed; no human sign-off; not a claim that RH is proved
% RH Notes blueprint source excerpts for Z1.1.03
% Verbatim excerpts in the notation of the parent blueprint. The labelled editions and source ranges appear below.
% See the unit YAML for the current proof status. Earlier-round passages retain their original wording.
\section*{Z1.1.03. Automorphy in the adjoint window}

% Parent source SHA-256: 9559fca8cc392b75d71f1c158fe0c635cfcab567e2f489b52db4c5bb29d98029
% blueprint/ch_P_r8_small.tex lines 419 through 437; r8:small:Z1.1.03
\paragraph{Z1.1.03. The initial local place and the gluing step.}
\phantomsection\label{r8:small:Z1.1.03}
In \nolinkurl{adjoint_automorphy_w_of_pure_weight_adjoint_lift}, the
level assumption makes \(M\) a multiple of \(\ell\).
Choose an embedding of the number field in \(\overline{\mathbb Q}_\ell\).
The prime defined by elements of norm \(<1\) contains \(\ell\) by S18,
hence contains \(M\). Apply the pointwise core comparison at this prime to
obtain a global representation \(\pi\) with the required data away from
\(M\). For any other place \(w\), use those data if \(w\nmid M\). If
\(w\mid M\), repeat the same pointwise comparison at \(w\) to obtain
the possibly different \(\pi_w\) and the ramified alternative there.
The gluing input accepts a witness separately for every \(w\), so this
choice supplies its full domain. The project-specific deductions still
required are weight-compatible descent and the pure-weight adjoint lift
which supply the core comparison; S08 is only one input to the former
(\nolinkurl{weight_compatible_descent}, \nolinkurl{adjoint_automorphy_w_of_printed_hida_window},
\path{UnitaryThreeInputsOfLiterature.lean:25--28};
\nolinkurl{adjoint_automorphy_w_of_pure_weight_adjoint_lift},
\path{DenseThresholdCases/MemOfUnder.lean:101--114}).

